\chapter*{ضمیمهٔ ۹: خانوادهٔ ایسم‌ها}
\addcontentsline{toc}{chapter}{ضمیمهٔ ۹: رابطهٔ ایسم‌ها با چپ و راست}

\begin{center}
\textit{تفکیک و توضیح ایدئولوژی‌های مرتبط}
\end{center}

\section*{خانوادهٔ چپ: تفاوت‌های درونی}

\begin{center}
\begin{tikzpicture}[scale=0.85, every node/.style={font=\small}]

% Root
\node[draw=leftred, fill=leftred!30, rounded corners, minimum width=3cm, font=\bfseries] (root) at (7,10) {\rl{سوسیالیسم}};

% First level branches
\node[draw=leftred!80, fill=leftred!20, rounded corners, minimum width=2.5cm] (marx) at (2,7) {\rl{مارکسیسم}};
\node[draw=leftred!60, fill=leftred!15, rounded corners, minimum width=2.5cm] (socdem) at (7,7) {\rl{سوسیال دموکراسی}};
\node[draw=leftred!80, fill=leftred!20, rounded corners, minimum width=2.5cm] (anar) at (12,7) {\rl{آنارشیسم}};

% Second level - Marxism branch
\node[draw=leftred!70, fill=leftred!10, rounded corners, minimum width=2cm] (lenin) at (0,4) {\rl{لنینیسم}};
\node[draw=leftred!70, fill=leftred!10, rounded corners, minimum width=2cm] (trot) at (2.5,4) {\rl{تروتسکیسم}};
\node[draw=leftred!70, fill=leftred!10, rounded corners, minimum width=2cm] (mao) at (5,4) {\rl{مائوئیسم}};

% Third level
\node[draw=leftred!60, fill=leftred!5, rounded corners] (stalin) at (0,1.5) {\rl{استالینیسم}};
\node[draw=leftred!60, fill=leftred!5, rounded corners] (euro) at (2.5,1.5) {\rl{یوروکمونیسم}};

% Social Democracy branch
\node[draw=centergreen!70, fill=centergreen!15, rounded corners, minimum width=2cm] (third) at (7,4) {\rl{راه سوم}};
\node[draw=centergreen!70, fill=centergreen!15, rounded corners, minimum width=2cm] (nordic) at (9.5,4) {\rl{مدل نوردیک}};

% Anarchism branch
\node[draw=violet!70, fill=violet!15, rounded corners, minimum width=2cm] (ancom) at (11,4) {\rl{آنارکو-کمونیسم}};
\node[draw=violet!70, fill=violet!15, rounded corners, minimum width=2cm] (ansyn) at (14,4) {\rl{آنارکو-سندیکالیسم}};

% Connections
\draw[->, line width=1.5pt] (root) -- (marx);
\draw[->, line width=1.5pt] (root) -- (socdem);
\draw[->, line width=1.5pt] (root) -- (anar);

\draw[->, line width=1pt] (marx) -- (lenin);
\draw[->, line width=1pt] (marx) -- (trot);
\draw[->, line width=1pt] (marx) -- (mao);

\draw[->, line width=1pt] (lenin) -- (stalin);
\draw[->, line width=1pt] (trot) -- (euro);

\draw[->, line width=1pt] (socdem) -- (third);
\draw[->, line width=1pt] (socdem) -- (nordic);

\draw[->, line width=1pt] (anar) -- (ancom);
\draw[->, line width=1pt] (anar) -- (ansyn);

\end{tikzpicture}
\end{center}

\subsection*{توضیح تفاوت‌های کلیدی در خانوادهٔ چپ}

\begin{tcolorbox}[
    enhanced,
    colback=leftred!5,
    colframe=leftred,
    boxrule=1.5pt,
    arc=5pt,
    fonttitle=\bfseries,
    title={\rl{🔴 سوسیالیسم در برابر کمونیسم}}
]

\textbf{سوسیالیسم:}
\begin{itemize}
    \item مالکیت اجتماعی بر ابزار تولید
    \item می‌تواند دموکراتیک باشد
    \item می‌تواند تدریجی باشد
    \item طیف گسترده‌ای از گرایش‌ها
\end{itemize}



\textbf{کمونیسم:}
\begin{itemize}
    \item جامعهٔ بی‌طبقه و بدون دولت
    \item مرحلهٔ نهایی تکامل تاریخ
    \item معمولاً انقلابی
    \item لغو کامل مالکیت خصوصی
\end{itemize}


\textbf{رابطه:} کمونیسم هدف نهایی است؛ سوسیالیسم مرحلهٔ گذار. اما در عمل، «سوسیالیسم» به رویکردهای معتدل‌تر و «کمونیسم» به رویکردهای انقلابی‌تر اطلاق می‌شود.
\end{tcolorbox}

\begin{tcolorbox}[
    enhanced,
    colback=leftred!5,
    colframe=leftred,
    boxrule=1.5pt,
    arc=5pt,
    fonttitle=\bfseries,
    title={\rl{🔴 مارکسیسم در برابر لنینیسم}}
]

\textbf{مارکسیسم (اصیل):}
\begin{itemize}
    \item انقلاب در کشورهای صنعتی پیشرفته
    \item پرولتاریای صنعتی موتور انقلاب
    \item انقلاب خودانگیخته
    \item دولت به تدریج محو می‌شود
\end{itemize}



\textbf{لنینیسم:}
\begin{itemize}
    \item انقلاب در «ضعیف‌ترین حلقهٔ زنجیر»
    \item حزب پیشاهنگ رهبری می‌کند
    \item انقلاب توطئه‌ای/سازماندهی‌شده
    \item دیکتاتوری پرولتاریا ضروری
\end{itemize}


\textbf{تفاوت کلیدی:} مارکس فکر می‌کرد انقلاب در آلمان یا انگلستان رخ می‌دهد. لنین «نظریه» را با واقعیت روسیه (کشور عقب‌مانده) تطبیق داد و نقش حزب را مرکزی کرد.
\end{tcolorbox}

\begin{tcolorbox}[
    enhanced,
    colback=leftred!5,
    colframe=leftred,
    boxrule=1.5pt,
    arc=5pt,
    fonttitle=\bfseries,
    title={\rl{🔴 تروتسکیسم در برابر استالینیسم}}
]

\textbf{تروتسکیسم:}
\begin{itemize}
    \item انقلاب مداوم/جهانی
    \item سوسیالیسم در یک کشور ناممکن
    \item دموکراسی درون حزبی
    \item ضد بوروکراسی
\end{itemize}



\textbf{استالینیسم:}
\begin{itemize}
    \item سوسیالیسم در یک کشور ممکن
    \item صنعتی‌سازی سریع به هر قیمت
    \item تمرکز قدرت در رهبر
    \item سرکوب مخالفان
\end{itemize}


\textbf{تاریخ:} تروتسکی و استالین پس از مرگ لنین (۱۹۲۴) رقابت کردند. استالین پیروز شد. تروتسکی تبعید و در ۱۹۴۰ ترور شد.
\end{tcolorbox}

\begin{tcolorbox}[
    enhanced,
    colback=centergreen!10,
    colframe=centergreen,
    boxrule=1.5pt,
    arc=5pt,
    fonttitle=\bfseries,
    title={\rl{🟢 سوسیال دموکراسی در برابر دموکراتیک سوسیالیسم}}
]

\textbf{سوسیال دموکراسی (امروز):}
\begin{itemize}
    \item پذیرش اقتصاد بازار
    \item اصلاح سرمایه‌داری
    \item دولت رفاه قوی
    \item تغییر تدریجی
\end{itemize}



\textbf{دموکراتیک سوسیالیسم:}
\begin{itemize}
    \item هدف: پایان سرمایه‌داری
    \item مالکیت اجتماعی واقعی
    \item از طریق دموکراسی
    \item برنی سندرز، کوربین
\end{itemize}


\textbf{تفاوت:} سوسیال دموکرات‌های امروز سرمایه‌داری را پذیرفته‌اند و فقط می‌خواهند آن را «اخلاقی‌تر» کنند. دموکراتیک سوسیالیست‌ها هنوز می‌خواهند از سرمایه‌داری فراتر بروند، اما از راه انتخابات.
\end{tcolorbox}

\section*{خانوادهٔ راست: تفاوت‌های درونی}

\begin{center}
\begin{tikzpicture}[scale=0.85, every node/.style={font=\small}]

% Root
\node[draw=rightblue, fill=rightblue!30, rounded corners, minimum width=3cm, font=\bfseries] (root) at (7,10) {\rl{راست}};

% First level branches
\node[draw=rightblue!80, fill=rightblue!20, rounded corners, minimum width=2.5cm] (conserv) at (2,7) {\rl{محافظه‌کاری}};
\node[draw=rightblue!60, fill=rightblue!15, rounded corners, minimum width=2.5cm] (liberal) at (7,7) {\rl{لیبرالیسم کلاسیک}};
\node[draw=rightblue!80, fill=rightblue!20, rounded corners, minimum width=2.5cm] (nation) at (12,7) {\rl{ناسیونالیسم}};

% Second level - Conservatism branch
\node[draw=rightblue!70, fill=rightblue!10, rounded corners, minimum width=2cm] (paleo) at (0,4) {\rl{پالئوکنسرواتیسم}};
\node[draw=rightblue!70, fill=rightblue!10, rounded corners, minimum width=2cm] (neo) at (3,4) {\rl{نئوکنسرواتیسم}};

% Liberalism branch
\node[draw=goldaccent!70, fill=goldaccent!15, rounded corners, minimum width=2cm] (neolib) at (6,4) {\rl{نئولیبرالیسم}};
\node[draw=goldaccent!70, fill=goldaccent!15, rounded corners, minimum width=2cm] (libert) at (9,4) {\rl{لیبرتارینیسم}};

% Nationalism branch
\node[draw=darkgray!70, fill=darkgray!15, rounded corners, minimum width=2cm] (civic) at (11,4) {\rl{ناسیونالیسم مدنی}};
\node[draw=darkgray!70, fill=darkgray!15, rounded corners, minimum width=2cm] (ethnic) at (14,4) {\rl{ناسیونالیسم قومی}};

% Third level
\node[draw=goldaccent!60, fill=goldaccent!5, rounded corners] (ancap) at (9,1.5) {\rl{آنارکو-کاپیتالیسم}};
\node[draw=darkgray!60, fill=darkgray!5, rounded corners] (fasc) at (14,1.5) {\rl{فاشیسم}};

% Connections
\draw[->, line width=1.5pt] (root) -- (conserv);
\draw[->, line width=1.5pt] (root) -- (liberal);
\draw[->, line width=1.5pt] (root) -- (nation);

\draw[->, line width=1pt] (conserv) -- (paleo);
\draw[->, line width=1pt] (conserv) -- (neo);

\draw[->, line width=1pt] (liberal) -- (neolib);
\draw[->, line width=1pt] (liberal) -- (libert);

\draw[->, line width=1pt] (nation) -- (civic);
\draw[->, line width=1pt] (nation) -- (ethnic);

\draw[->, line width=1pt] (libert) -- (ancap);
\draw[->, line width=1pt] (ethnic) -- (fasc);

\end{tikzpicture}
\end{center}

\subsection*{توضیح تفاوت‌های کلیدی در خانوادهٔ راست}

\begin{tcolorbox}[
    enhanced,
    colback=rightblue!5,
    colframe=rightblue,
    boxrule=1.5pt,
    arc=5pt,
    fonttitle=\bfseries,
    title={\rl{🔵 محافظه‌کاری در برابر لیبرالیسم کلاسیک}}
]

\textbf{محافظه‌کاری:}
\begin{itemize}
    \item حفظ سنت و نهادها
    \item شک به تغییر سریع
    \item اهمیت جامعه و خانواده
    \item گاهی دولت‌گرا (در مسائل اخلاقی)
    \item برک، راسل کرک
\end{itemize}



\textbf{لیبرالیسم کلاسیک:}
\begin{itemize}
    \item آزادی فردی محور
    \item حقوق طبیعی
    \item بازار آزاد
    \item دولت محدود
    \item لاک، اسمیت، میل
\end{itemize}


\textbf{رابطه:} در آمریکا این دو ترکیب شده‌اند («فیوژنیسم»). اما تنش همیشه هست: آیا دولت باید «ارزش‌ها» را ترویج کند (محافظه‌کار) یا بی‌طرف باشد (لیبرال)؟
\end{tcolorbox}

\begin{tcolorbox}[
    enhanced,
    colback=rightblue!5,
    colframe=rightblue,
    boxrule=1.5pt,
    arc=5pt,
    fonttitle=\bfseries,
    title={\rl{🔵 نئولیبرالیسم در برابر لیبرتارینیسم}}
]

\textbf{نئولیبرالیسم:}
\begin{itemize}
    \item بازار آزاد + دولت فعال
    \item خصوصی‌سازی، مقررات‌زدایی
    \item جهانی‌سازی
    \item پذیرش دولت (برای حمایت از بازار)
    \item تاچر، ریگان
\end{itemize}



\textbf{لیبرتارینیسم:}
\begin{itemize}
    \item بازار آزاد + حداقل دولت
    \item ضد همهٔ دخالت‌ها
    \item ممکن است ضد جهانی‌سازی باشد
    \item دولت = دشمن
    \item راتبارد، ران پال
\end{itemize}


\textbf{تفاوت کلیدی:} نئولیبرال‌ها دولت را لازم می‌دانند (برای حفظ بازار، ارتش، و...). لیبرتارین‌ها می‌خواهند دولت حداقلی یا صفر باشد.
\end{tcolorbox}

\begin{tcolorbox}[
    enhanced,
    colback=rightblue!5,
    colframe=rightblue,
    boxrule=1.5pt,
    arc=5pt,
    fonttitle=\bfseries,
    title={\rl{🔵 پالئوکنسرواتیسم در برابر نئوکنسرواتیسم}}
]

\textbf{پالئوکنسرواتیسم:}
\begin{itemize}
    \item ناسیونالیسم اقتصادی
    \item حمایت‌گرایی تجاری
    \item محدودیت مهاجرت
    \item انزواطلبی در سیاست خارجی
    \item پت بوکانان
\end{itemize}



\textbf{نئوکنسرواتیسم:}
\begin{itemize}
    \item تجارت آزاد جهانی
    \item مهاجرت نسبتاً باز
    \item مداخله‌جویی خارجی
    \item صدور دموکراسی
    \item کریستول، بوش
\end{itemize}


\textbf{تاریخ:} نئوکنسرواتیسم از چپ آمد (یهودیان نیویورکی که از چپ مأیوس شدند). پالئوکنسرواتیسم «محافظه‌کاری قدیمی» آمریکاست. ترامپ به پالئوکنسرواتیسم نزدیک‌تر است.
\end{tcolorbox}

\section*{ایسم‌های دیگر و رابطه‌شان با چپ و راست}

\begin{table}[H]
\centering
\small
\begin{tabular}{|>{\columncolor{darkgray!10}}r|c|p{7cm}|}
\hline
\rowcolor{darkgray!30}
\textbf{ایسم} & \textbf{جایگاه در طیف} & \textbf{توضیح} \\
\hline
\textbf{فمینیسم} & عمدتاً چپ & برابری جنسیتی؛ اما فمینیسم لیبرال و رادیکال متفاوت‌اند \\
\hline
\textbf{محیط‌زیست‌گرایی} & عمدتاً چپ & جنبش سبز؛ اما «محافظه‌کاری سبز» هم وجود دارد \\
\hline
\textbf{پوپولیسم} & هر دو & سبک سیاسی، نه محتوا؛ پوپولیسم چپ و راست داریم \\
\hline
\textbf{ناسیونالیسم} & عمدتاً راست (امروز) & تاریخاً هم چپ داشته (جنبش‌های آزادی‌بخش) \\
\hline
\textbf{دین‌سالاری} & راست افراطی & حکومت دینی؛ تئوکراسی \\
\hline
\textbf{سکولاریسم} & عمدتاً چپ/میانه & جدایی دین و دولت \\
\hline
\textbf{کمونیتارینیسم} & فراتر از طیف & نقد فردگرایی لیبرال؛ هم از چپ هم از راست طرفدار دارد \\
\hline
\textbf{فدرالیسم} & خنثی & ساختار حکومتی؛ نه چپ نه راست ذاتاً \\
\hline
\textbf{رپوبلیکانیسم} & خنثی تا چپ & تأکید بر شهروندی و مشارکت؛ جمهوری‌خواهی \\
\hline
\textbf{پراگماتیسم} & خنثی & روش، نه محتوا؛ «هر چه کار کند» \\
\hline
\textbf{اگزیستانسیالیسم} & معمولاً چپ & سارتر چپ بود؛ اما هایدگر نازی شد \\
\hline
\textbf{پست‌مدرنیسم} & چپ فرهنگی & نقد روایت‌های بزرگ؛ راست آن را دشمن می‌داند \\
\hline
\textbf{کاپیتالیسم} & راست & مالکیت خصوصی و بازار آزاد \\
\hline
\textbf{کورپوراتیسم} & راست اقتدارگرا & سازماندهی صنفی؛ فاشیسم ایتالیا \\
\hline
\textbf{سندیکالیسم} & چپ & قدرت به اتحادیه‌های کارگری \\
\hline
\textbf{پلورالیسم} & لیبرال/میانه & پذیرش تنوع گروه‌ها و دیدگاه‌ها \\
\hline
\end{tabular}
\caption{ایسم‌های مختلف و جایگاه‌شان در طیف سیاسی}
\end{table}

\section*{نمودار جامع روابط}

\begin{center}
\begin{tikzpicture}[scale=0.7, every node/.style={font=\scriptsize}]

% Main axis
\draw[line width=3pt, left color=leftred, right color=rightblue] (0,0) -- (18,0);
\node[below, font=\bfseries] at (0,0) {\rl{چپ افراطی}};
\node[below, font=\bfseries] at (9,0) {\rl{میانه}};
\node[below, font=\bfseries] at (18,0) {\rl{راست افراطی}};

% Ideologies positioned
% Far Left
\node[draw=leftred, fill=leftred!30, rounded corners, rotate=45] at (1,1.5) {\rl{کمونیسم}};
\node[draw=leftred, fill=leftred!25, rounded corners, rotate=45] at (2,2) {\rl{آنارکو-کمونیسم}};

% Left
\node[draw=leftred!80, fill=leftred!20, rounded corners, rotate=45] at (3.5,1.5) {\rl{مارکسیسم}};
\node[draw=leftred!70, fill=leftred!15, rounded corners, rotate=45] at (4.5,2.5) {\rl{سوسیالیسم}};
\node[draw=leftred!60, fill=leftred!10, rounded corners, rotate=45] at (5.5,1.8) {\rl{چپ نو}};

% Center-Left
\node[draw=centergreen!80, fill=centergreen!20, rounded corners, rotate=45] at (7,2) {\rl{سوسیال دموکراسی}};
\node[draw=centergreen!70, fill=centergreen!15, rounded corners, rotate=45] at (8,1.5) {\rl{لیبرالیسم اجتماعی}};

% Center
\node[draw=centergreen, fill=centergreen!30, rounded corners] at (9,3) {\rl{میانه‌روی}};
\node[draw=violet!70, fill=violet!15, rounded corners, rotate=45] at (9.5,1.8) {\rl{راه سوم}};

% Center-Right
\node[draw=rightblue!60, fill=rightblue!10, rounded corners, rotate=45] at (10.5,2) {\rl{لیبرالیسم کلاسیک}};
\node[draw=rightblue!70, fill=rightblue!15, rounded corners, rotate=45] at (11.5,1.5) {\rl{دموکراسی مسیحی}};

% Right
\node[draw=rightblue!80, fill=rightblue!20, rounded corners, rotate=45] at (13,2) {\rl{محافظه‌کاری}};
\node[draw=goldaccent!70, fill=goldaccent!15, rounded corners, rotate=45] at (14,2.5) {\rl{نئولیبرالیسم}};
\node[draw=goldaccent!80, fill=goldaccent!20, rounded corners, rotate=45] at (15,1.8) {\rl{لیبرتارینیسم}};

% Far Right
\node[draw=rightblue, fill=rightblue!30, rounded corners, rotate=45] at (16.5,2) {\rl{ناسیونالیسم افراطی}};
\node[draw=darkgray, fill=darkgray!30, rounded corners, rotate=45] at (17.5,1.5) {\rl{فاشیسم}};

% Vertical axis label
\node[left, rotate=90, font=\bfseries] at (-0.5,2) {\rl{اقتدارگرایی ← آزادی‌خواهی}};

\end{tikzpicture}
\end{center}

\section*{واژه‌های گیج‌کننده: روشن‌سازی}

\begin{tcolorbox}[
    enhanced,
    colback=goldaccent!10,
    colframe=goldaccent,
    boxrule=1.5pt,
    arc=5pt,
    fonttitle=\bfseries,
    title={\rl{🟡 لیبرال: دو معنای متفاوت}}
]

\textbf{لیبرال در آمریکا:}
\begin{itemize}
    \item = چپ‌میانه
    \item طرفدار دولت رفاه
    \item حقوق اقلیت‌ها
    \item دموکرات‌ها
    \item مثال: باراک اوباما
\end{itemize}



\textbf{لیبرال در اروپا/جهان:}
\begin{itemize}
    \item = طرفدار بازار آزاد
    \item ضد دخالت دولت
    \item آزادی اقتصادی
    \item گاه راست‌میانه
    \item مثال: FDP آلمان
\end{itemize}


\textbf{چرا این سردرگمی؟} در آمریکا، «لیبرال» از دههٔ ۱۹۳۰ (نیو دیل) معنای چپ‌میانه گرفت. در اروپا، «لیبرال» همچنان به معنای کلاسیک (آزادی از دولت) باقی ماند.
\end{tcolorbox}

\begin{tcolorbox}[
    enhanced,
    colback=goldaccent!10,
    colframe=goldaccent,
    boxrule=1.5pt,
    arc=5pt,
    fonttitle=\bfseries,
    title={\rl{🟡 جمهوری‌خواه: دو معنا}}
]

\textbf{جمهوری‌خواهی فلسفی:}
\begin{itemize}
    \item ضد سلطنت
    \item تأکید بر شهروندی
    \item فضیلت مدنی
    \item مشارکت عمومی
    \item سنت: روم، ماکیاولی
\end{itemize}



\textbf{حزب جمهوری‌خواه آمریکا:}
\begin{itemize}
    \item یکی از دو حزب بزرگ
    \item امروز: راست/محافظه‌کار
    \item تاریخاً: ضد برده‌داری (لینکلن)
    \item ربطی به جمهوری‌خواهی فلسفی ندارد!
\end{itemize}

\end{tcolorbox}

\begin{tcolorbox}[
    enhanced,
    colback=goldaccent!10,
    colframe=goldaccent,
    boxrule=1.5pt,
    arc=5pt,
    fonttitle=\bfseries,
    title={\rl{🟡 دموکرات: گیجی‌های مشابه}}
]
\begin{itemize}
    \item \textbf{دموکراسی:} حکومت مردم (فلسفی)
    \item \textbf{حزب دموکرات آمریکا:} چپ‌میانه (امروز)
    \item \textbf{دموکرات‌های مسیحی:} راست‌میانهٔ اروپا
    \item \textbf{دموکرات‌های اجتماعی:} = سوسیال دموکرات‌ها
    \item \textbf{لیبرال دموکرات‌ها:} میانهٔ بریتانیا
\end{itemize}
\end{tcolorbox}

\section*{خلاصه: خانواده‌های ایدئولوژیک}

\begin{table}[H]
\centering
\begin{tabular}{|>{\columncolor{leftred!15}}r|p{10cm}|}
\hline
\rowcolor{leftred!30}
\multicolumn{2}{|c|}{\textbf{خانوادهٔ چپ}} \\
\hline
\textbf{ریشه مشترک} & نقد نابرابری؛ تأکید بر عدالت اجتماعی \\
\hline
\textbf{اعضا} & سوسیالیسم، کمونیسم، مارکسیسم، سوسیال دموکراسی، آنارشیسم چپ \\
\hline
\textbf{اختلاف اصلی} & انقلاب یا اصلاح؟ دولت یا بی‌دولت؟ \\
\hline
\end{tabular}

\vspace{5mm}

\begin{tabular}{|>{\columncolor{rightblue!15}}r|p{10cm}|}
\hline
\rowcolor{rightblue!30}
\multicolumn{2}{|c|}{\textbf{خانوادهٔ راست}} \\
\hline
\textbf{ریشه مشترک} & دفاع از سنت، مالکیت، یا آزادی اقتصادی \\
\hline
\textbf{اعضا} & محافظه‌کاری، لیبرالیسم کلاسیک، نئولیبرالیسم، لیبرتارینیسم، ناسیونالیسم \\
\hline
\textbf{اختلاف اصلی} & سنت یا بازار؟ ملت یا جهان؟ \\
\hline
\end{tabular}

\vspace{5mm}

\begin{tabular}{|>{\columncolor{centergreen!15}}r|p{10cm}|}
\hline
\rowcolor{centergreen!30}
\multicolumn{2}{|c|}{\textbf{خانوادهٔ میانه}} \\
\hline
\textbf{ریشه مشترک} & پراگماتیسم؛ رد افراط؛ ترکیب عناصر \\
\hline
\textbf{اعضا} & لیبرالیسم اجتماعی، راه سوم، دموکراسی مسیحی \\
\hline
\textbf{اختلاف اصلی} & تا کجا به چپ؟ تا کجا به راست؟ \\
\hline
\end{tabular}
\end{table}

\vspace{10mm}
\begin{center}
\rule{0.6\textwidth}{1pt}

\vspace{3mm}
\textit{«همهٔ ایدئولوژی‌ها یک هستهٔ حقیقت دارند—و یک حاشیهٔ اشتباه.»}

\vspace{3mm}
\rule{0.6\textwidth}{1pt}
\end{center}